\chapter*{Governance}
\addcontentsline{toc}{chapter}{Governance}

Three students, one supervisor, no fixed roles. Figure~\ref{fig:governance} shows how the team is
organised. The supervisor is at the top, because he sets the requirements and reviews the work; the
dashed arrows are the questions on the requirements that each student can bring to him. The three
students are on the same level, joined by the lines of their joint decisions, and each of them
works in every phase. The arrows at the bottom converge on the shared repository, where all the work
ends up.

% The team: supervisor, three members on the same level, one shared repository
\begin{figure}[H]
\centering
  \begin{tikzpicture}[font=\footnotesize,
      person/.style={rounded corners=4pt, cat infra, text width=3.6cm, align=center,
                     minimum height=1.15cm, inner sep=4pt},
      link/.style={-{Stealth[length=2mm]}, black!55, line width=0.6pt}]
    \node[rounded corners=4pt, draw=accent, fill=accent!8, text width=4.6cm, align=center,
          inner sep=5pt] (sup) at (7.9,0)
      {{\color{accent}\Large\faIcon{chalkboard-teacher}}\\[2pt]\textbf{Prof.\ Flavio Corradini}\\
       {\scriptsize\color{black!65}supervisor: requirements, reviews}};
    \node[person] (m1) at (1.9,-2.6) {{\color{accent}\faUser}\enspace\textbf{Matteo Rizzuto}\\
      {\scriptsize\color{black!65}every phase}};
    \node[person] (m2) at (7.9,-2.6) {{\color{accent}\faUser}\enspace\textbf{Mattia Farabollini}\\
      {\scriptsize\color{black!65}every phase}};
    \node[person] (m3) at (13.9,-2.6) {{\color{accent}\faUser}\enspace\textbf{Paolo Campanari}\\
      {\scriptsize\color{black!65}every phase}};
    \node[rounded corners=4pt, cat data, text width=5.2cm, align=center, inner sep=5pt]
      (repo) at (7.9,-4.9) {{\Large\faGithub}\\[2pt]\textbf{Shared Git repository}\\
      {\scriptsize\color{black!65}code and report, visible to all}};
    % Questions on the requirements go to the supervisor
    \foreach \m in {m1, m2, m3}
      \draw[link, dashed] (\m.north) -- (sup.south);
    \node[fill=white, inner sep=1.5pt, font=\scriptsize\itshape, text=black!65] at (7.9,-1.35)
      {questions on the requirements};
    % A flat team: joint decisions
    \draw[black!40, line width=1.2pt] (m1.east) -- (m2.west);
    \draw[black!40, line width=1.2pt] (m2.east) -- (m3.west);
    \node[above, font=\scriptsize\itshape, text=black!65, align=center] at (4.9,-2.6)
      {joint\\decisions};
    \node[above, font=\scriptsize\itshape, text=black!65, align=center] at (10.9,-2.6)
      {joint\\decisions};
    % Everyone works on the repository
    \foreach \m in {m1, m2, m3}
      \draw[link] (\m.south) -- (repo.north);
  \end{tikzpicture}
  \caption{The organisation of the team: one supervisor, three members on the same level and one
    shared repository.}
  \label{fig:governance}
\end{figure}

The team has no fixed roles. All three students take part in every phase of the project:
domain analysis, architecture, development, evaluation and the writing of this report. The reason
is resilience: since everyone knows the whole system, the project never depends on a single
person, and anyone can continue the work of the others when needed.

Tasks are not assigned once and for all. They are distributed according to everyone's
availability and to the priorities of the Gantt diagram (Figure~\ref{fig:gantt}), and anyone can
take over a task that someone else has started. The decisions that shape the project, such as its
scope, its architecture and the way it is evaluated, are taken together by the three members.
Questions on the requirements are brought to the supervisor, Prof.\ Flavio Corradini, who also
reviews the progress of the work.

All the work lives in one shared Git repository, which holds both the code and this report and is
visible to everyone, so that each member always sees the current state of the project.
